\section*{Sets and Indices}

\[
P=\{A,B\}
\]
set of products, indexed by \(p\).

\section*{Parameters}

All parameter names below match the data fields.

For each product \(p \in P\):
\begin{align*}
s_p &: \text{steel required per unit (from \texttt{Steel\_Required\_kg})},\\
a_p &: \text{aluminum required per unit (from \texttt{Aluminum\_Required\_kg})},\\
\ell_p &: \text{labor required per unit (from \texttt{Labor\_Required\_hours})},\\
r_p &: \text{profit per unit (from \texttt{Profit\_yuan})},\\
f_p &: \text{fixed setup cost (from \texttt{Setup\_Cost\_yuan})}.
\end{align*}

Resource and policy parameters:
\begin{align*}
S &:= \texttt{available\_steel},\\
A &:= \texttt{available\_aluminum},\\
L &:= \texttt{available\_labor},\\
c^{ot} &:= \texttt{overtime\_pay},\\
\overline{O} &:= \texttt{max\_overtime},\\
T^{rev} &:= \texttt{overtime\_review\_threshold},\\
F^{rev} &:= \texttt{overtime\_review\_fee},\\
T^A &:= \texttt{product\_A\_line\_audit\_threshold},\\
F^A &:= \texttt{product\_A\_line\_audit\_fee}.
\end{align*}

Implementation constants used in the code:
\[
M = 10000, \qquad M^{rev} = \texttt{max\_overtime}.
\]

\section*{Decision Variables}

\begin{align*}
x_p &\ge 0 \qquad &&\text{production quantity of product } p \text{ (code: \texttt{x[p]})},\\
y_p &\in \{0,1\} \qquad &&\text{1 if product } p \text{ is produced/setup is paid (code: \texttt{y[p]})},\\
o &\in [0,\overline{O}] \qquad &&\text{overtime hours (code: \texttt{overtime})},\\
z^{rev} &\in \{0,1\} \qquad &&\text{1 if overtime review fee is triggered (code: \texttt{review})},\\
z^{A} &\in \{0,1\} \qquad &&\text{1 if Product A line-audit fee is triggered (code: \texttt{product\_A\_line\_audit})}.
\end{align*}

\section*{Objective}

Maximize net profit:
\[
\max \;
\sum_{p \in P} r_p x_p
-\sum_{p \in P} f_p y_p
-c^{ot} o
-F^{rev} z^{rev}
-F^{A} z^{A}.
\]

\section*{Constraints}

Resource limits:
\begin{align*}
\sum_{p \in P} s_p x_p &\le S,\\
\sum_{p \in P} a_p x_p &\le A,\\
\sum_{p \in P} \ell_p x_p &\le L + o.
\end{align*}

Setup-linking constraints:
\[
x_p \le M y_p \qquad \forall p \in P.
\]

Overtime review gate:
\[
o \le T^{rev} + M^{rev} z^{rev}.
\]

Product A line-audit gate:
\[
x_A \le T^A + M z^{A}.
\]

Variable domains:
\[
x_p \ge 0 \;\; (\forall p \in P), \qquad
0 \le o \le \overline{O}, \qquad
y_p \in \{0,1\} \;\; (\forall p \in P), \qquad
z^{rev}, z^{A} \in \{0,1\}.
\]

\section*{Notes on Implementation}

The formulation above matches the implemented mixed-integer linear model in \texttt{current\_heuristic.py}.

The overtime-review rule is modeled as a binary gate, not as a proportional surcharge: if \(o \le T^{rev}\), then \(z^{rev}=0\) is feasible and no review fee is paid; if \(o>T^{rev}\), feasibility requires \(z^{rev}=1\), which adds the one-time fee \(F^{rev}\).

Likewise, the Product A audit rule is modeled as a binary gate: up to \(\texttt{product\_A\_line\_audit\_threshold}\) units of A can be produced with \(z^{A}=0\); producing more forces \(z^{A}=1\), which adds the one-time fee \(\texttt{product\_A\_line\_audit\_fee}\). This audit fee is separate from the overtime review fee.

The code uses the big-\(M\) value \(10000\) both for setup linking \(x_p \le 10000\,y_p\) and for the Product A audit gate \(x_A \le T^A + 10000\,z^A\). The overtime review gate uses \(\texttt{max\_overtime}\) as its big-\(M\) term exactly as written in the source code.

The model is solved as a MILP through \texttt{GUROBI\_CMD(msg=0)}.
